% Volume IV: Law as an Epistemic Machine
% Part VII. Evidence Before Action
% Chapter 38: Cross-Examination as Adversarial Testing
% Spine: proof-spines/ch038-spine.md

\chapter{Cross-Examination as Adversarial Testing}
\label{ch:ch038}

Cross-examination probes whether a claim survives transformations of perspective. A witness's assertion is not simply counted as one more datum; its dependence on memory, incentives, vantage, language, and prior interpretation is made contestable in public. The chapter captures this as a \textbf{procedural robustness} functional:
\[
R(H)
=
\inf_{\delta\in\Delta}
P(H\mid E+\delta),
\]
\Definition\ where \(\Delta\) contains the admissible perturbations or challenges generated by questioning.
A conclusion that collapses under minor, procedurally allowed challenges has low robustness even if it looked persuasive before testing. Cross-examination thus measures not bare confidence, but the stability of confidence under organized pressure.
